\subsection{Targets covered by recent vanishing theorems}
\begin{proposition}[A1Q: quasi-F-split and Fano targets; Known deduction]
\label{a:known-targets}
Let $X$ be a smooth connected projective $n$-fold over an algebraically
closed field of characteristic $p>0$, and let $Y$ be a smooth effective
ample Cartier divisor. The following statements hold for this specified
divisor, and hence for a smooth hyperplane of a specified embedding.
\begin{enumerate}
\item If $Y$ is quasi-F-split and $n\geq3$, then
$\Pic(X)\to\Pic(Y)$ is injective. If also $n\geq4$, it is an isomorphism.
\item If $n=4$ and $Y$ is Fano, then
$\Pic(X)\to\Pic(Y)$ is an isomorphism, including in characteristic two.
\end{enumerate}
Here quasi-F-split means that, for some integer $r\geq1$, the natural
$\OO_{Y^{(1)}}$-linear map
$\OO_{Y^{(1)}}\to F_*W_r(\OO_Y)/p$ has an
$\OO_{Y^{(1)}}$-linear retraction, with the Frobenius module structure
as in \cite[\S1, p.~1]{RootPetrov2025}.
\end{proposition}
\begin{proof}
Put $d=\dim Y=n-1$ and $A=\OO_Y(Y)$, which is ample.
Petrov's Corollary~4.10(2), printed p.~16 of the opened version,
gives
\[
 H^{d-i}(Y,\omega_Y\otimes A^m)=0
       \qquad (m>0,\ 0<i<d)
\]
in the first case: use differential-form degree $d$, cohomological
degree $d-i$, and ample bundle $A^m$. The sum of the two degrees is
$2d-i>d$. Serre duality identifies the dual of the displayed group
with $H^i(Y,A^{-m})$. Thus all the negative first cohomology groups
vanish for $n\geq3$, and all the negative second groups vanish for
$n\geq4$. Apply Corollary~\ref{a:vanishing}.

For the second case, Kawakami--Tanaka prove Akizuki--Nakano vanishing
for every smooth Fano threefold in every positive characteristic:
\cite[Theorem~7.2, p.~56]{RootKawakamiTanakaII2025}, together with
\cite[Theorem~6.3, p.~50]{RootKawakamiTanakaI2026} for the case
$\Pic(Y)=\mathbf ZK_Y$ with very ample anticanonical bundle.
Take differential-form degree zero, cohomological degree $i=1,2$,
and ample divisor corresponding to $A^m$. Both sums of degrees are
strictly less than three, so $H^i(Y,A^{-m})=0$ for every $m>0$.
Corollary~\ref{a:vanishing} again applies. Neither proof requires
that the ambient variety lift, that the given divisor lift, or that
$p$ exceed its dimension.
\end{proof}

These are target hypotheses. They are not properties established for
an arbitrary smooth ample divisor. In particular this proposition
does not settle A1 or A2. The cited papers also distinguish Fano
threefolds from quasi-F-split varieties: some smooth Fano threefolds
are not quasi-F-split, so the second assertion is not redundant.
